% =============================================================
%  1. OBJETIVO, EQUIPAMIENTO Y MONTAJE (común)
% =============================================================
\renewcommand{\seccorta}{Equipamiento}
\section{Objetivo y equipamiento}

\begin{frame}{Objetivo y metodología}
  \begin{block}{Parte 2}
    Llevar a la práctica los conceptos de la Parte 1 con un enlace real de corta distancia
    entre dos LiteBeam 5AC Gen2: configuración, análisis de espectro, alineación y pruebas de
    rendimiento.
  \end{block}
  \vspace{4pt}
  \begin{columns}[T]
    \column{0.48\textwidth}
    \textbf{Etapas}
    \begin{enumerate}
      \item Reconocimiento del equipo y montaje.
      \item Configuración básica en el laboratorio.
      \item Pruebas en el pasillo de la facultad (70 m).
    \end{enumerate}
    \column{0.48\textwidth}
    \textbf{Roles}
    \begin{itemize}
      \item \textbf{Punto A} (AP PtP, ``transmisor''): \grupoAP
      \item \textbf{Punto B} (Station PtP, ``receptor''): \grupoST
    \end{itemize}
  \end{columns}
\end{frame}

\begin{frame}{Ubiquiti LiteBeam 5AC Gen2}
  \begin{columns}[c]
    \column{0.42\textwidth}\centering
    \figura{litebeam.jpg}{\linewidth}
    \column{0.55\textwidth}
    \small
    \begin{itemize}
      \item \textbf{Alimentador (InnerFeed):} radio, antena, alimentación y red integrados; sin cables coaxiales.
      \item \textbf{Reflector parabólico:} 358 $\times$ 272 mm, perforado para reducir la carga de viento.
      \item \textbf{Montaje:} ajuste independiente de azimut y elevación, con nivel de burbuja.
      \item \textbf{Puerto RJ45 Gigabit:} datos + PoE pasivo 24 V / 0.3 A.
    \end{itemize}
  \end{columns}
\end{frame}

\begin{frame}{Partes del equipo}
  \begin{columns}[c]
    \column{0.32\textwidth}\centering
    \figura{vista_inferior.jpg}{\linewidth}\\
    {\scriptsize Vista inferior}
    \column{0.32\textwidth}\centering
    \figura{vista_lateral.jpg}{\linewidth}\\
    {\scriptsize Vista lateral}
    \column{0.32\textwidth}
    \small
    \textcolor[HTML]{7CB342}{\textbf{Verde:}} alimentador con puerto PoE y LEDs\\[3pt]
    \textcolor[HTML]{F9A825}{\textbf{Amarillo:}} puerto RJ45\\[3pt]
    \textcolor[HTML]{E53935}{\textbf{Rojo:}} ajuste de azimut\\[3pt]
    \textcolor[HTML]{1E88E5}{\textbf{Azul:}} ajuste de elevación\\[3pt]
    \textcolor[HTML]{7E57C2}{\textbf{Violeta:}} nivel de burbuja
  \end{columns}
\end{frame}

\begin{frame}{Montaje y conexión mediante inyector PoE}
  \begin{columns}[c]
    \column{0.3\textwidth}\centering
    \figura{montaje.jpg}{\linewidth}\\
    {\scriptsize Equipos preensamblados sobre trípodes}
    \column{0.34\textwidth}\centering
    \figura{poe.jpg}{\linewidth}
    \column{0.32\textwidth}
    \small
    \begin{enumerate}
      \item Solo alimentación.
      \item Alimentación + LAN hacia PC, switch o router.
      \item Gestión inalámbrica desde el celular (radio Wi-Fi de 2.4 GHz).
    \end{enumerate}
    \vspace{4pt}
    \textbf{Laboratorio:} ambos puntos cableados a una PC.\\
    \textbf{Pasillo:} Punto B solo alimentado y gestionado por Wi-Fi.
  \end{columns}
\end{frame}
